\documentclass[10pt,aspectratio=169]{beamer}

% ------------------------------------------------------------------
% Presentació adaptada: Probabilitat 2 amb orientació computacional
% Grau en Enginyeria Informàtica - Estadística
% ------------------------------------------------------------------

\usepackage[utf8]{inputenc}
\usepackage[T1]{fontenc}
\usepackage[catalan]{babel}
\usepackage{amsmath,amssymb}
\usepackage{booktabs}
\usepackage{array}
\usepackage{tikz}
\usepackage{pgfplots}
\usepackage{xcolor}
\usepackage{listings}
\usetikzlibrary{positioning,arrows.meta,calc,shapes.geometric}
\pgfplotsset{compat=1.18}

% ------------------------------------------------------------------
% Estil visual coherent amb les presentacions anteriors
% ------------------------------------------------------------------
\definecolor{UABGreen}{RGB}{0,103,71}
\definecolor{UABLight}{RGB}{224,241,236}
\definecolor{DeepBlue}{RGB}{20,65,110}
\definecolor{AccentOrange}{RGB}{230,126,34}
\definecolor{SoftGray}{RGB}{245,247,250}
\definecolor{CodeBack}{RGB}{248,248,248}
\definecolor{CodeFrame}{RGB}{210,215,220}

\usetheme{Madrid}
\usecolortheme[named=UABGreen]{structure}
\setbeamercolor{title}{fg=white,bg=UABGreen}
\setbeamercolor{frametitle}{fg=white,bg=UABGreen}
\setbeamercolor{block title}{fg=white,bg=DeepBlue}
\setbeamercolor{block body}{fg=black,bg=SoftGray}
\setbeamercolor{alerted text}{fg=AccentOrange}
\setbeamertemplate{navigation symbols}{}
\setbeamertemplate{itemize item}{\color{UABGreen}$\blacktriangleright$}
\setbeamertemplate{itemize subitem}{\color{AccentOrange}$\bullet$}
\setbeamertemplate{enumerate item}{\color{UABGreen}\insertenumlabel.}

\setbeamertemplate{footline}{%
  \leavevmode%
  \hbox{%
  \begin{beamercolorbox}[wd=.72\paperwidth,ht=2.4ex,dp=1ex,leftskip=1.4em]{author in head/foot}%
    Estadística -- Grau en Enginyeria Informàtica
  \end{beamercolorbox}%
  \begin{beamercolorbox}[wd=.28\paperwidth,ht=2.4ex,dp=1ex,center]{date in head/foot}%
    \insertframenumber{} / \inserttotalframenumber
  \end{beamercolorbox}}%
  \vskip0pt%
}

\lstdefinelanguage{R}{%
  morekeywords={if,else,repeat,while,function,for,in,next,break,TRUE,FALSE,NULL,NA,NaN,Inf,
    library,mean,median,sd,var,summary,table,prop.table,xtabs,addmargins,round,
    ggplot,aes,geom_point,geom_line,geom_col,geom_tile,geom_smooth,facet_wrap,
    labs,theme_minimal,summarise,group_by,mutate,read_csv,filter,arrange,set.seed,
    rbinom,rgeom,rpois,dbinom,dgeom,dpois,pbinom,ppois,runif,sample,replicate,
    cumsum,cummean,tibble,rep,seq,case_when,future_map_dbl,plan,multisession,
    furrr,pmap_dbl,rowwise,data.frame},
  sensitive=true,
  morecomment=[l]\#,
  morestring=[b]",
  morestring=[b]'
}
\lstset{%
  language=R,
  basicstyle=\ttfamily\scriptsize,
  keywordstyle=\color{DeepBlue}\bfseries,
  commentstyle=\color{gray},
  stringstyle=\color{AccentOrange},
  backgroundcolor=\color{CodeBack},
  frame=single,
  rulecolor=\color{CodeFrame},
  breaklines=true,
  columns=fullflexible,
  keepspaces=true,
  showstringspaces=false,
  literate={à}{{\`a}}1 {è}{{\`e}}1 {é}{{\'e}}1 {í}{{\'i}}1 {ï}{{\"i}}1 {ò}{{\`o}}1 {ó}{{\'o}}1 {ú}{{\'u}}1 {ü}{{\"u}}1 {ç}{{\c{c}}}1 {À}{{\`A}}1 {È}{{\`E}}1 {É}{{\'E}}1 {Í}{{\'I}}1 {Ò}{{\`O}}1 {Ó}{{\'O}}1 {Ú}{{\'U}}1
}

\newcommand{\R}{\textsf{R}}
\newcommand{\RStudio}{\textsf{RStudio/Posit}}
\newcommand{\GitHub}{\textsf{GitHub}}
\newcommand{\important}[1]{\textcolor{UABGreen}{\textbf{#1}}}
\newcommand{\exemple}[1]{\textcolor{AccentOrange}{\textbf{#1}}}
\newcommand{\Prob}{\mathbb{P}}
\newcommand{\E}{\mathbb{E}}
\newcommand{\Var}{\operatorname{Var}}

% ------------------------------------------------------------------
% Dades de la presentació
% ------------------------------------------------------------------
\title[Variables aleatòries discretes]{Probabilitat II}
\subtitle{Variables aleatòries discretes, esperança, variància i distribucions computacionals}
\author{Estadística -- Grau en Enginyeria Informàtica}
\date{Curs 2025--26}

\begin{document}

% ------------------------------------------------------------------
\begin{frame}[plain]
  \titlepage
  \vspace{-0.45cm}
  \begin{center}
  \small
  \textcolor{DeepBlue}{De comptar esdeveniments a simular sistemes: errors, intents, connexions i càrrega}
  \end{center}
\end{frame}

% ------------------------------------------------------------------
\begin{frame}{Per què variables aleatòries discretes en informàtica?}
\begin{columns}[T]
\column{0.55\textwidth}
\begin{itemize}
  \item Molts fenòmens informàtics són \important{comptatges}:
  \begin{itemize}
    \item nombre de paquets perduts,
    \item nombre d'errors en una bateria de tests,
    \item nombre d'intents fins al primer èxit,
    \item nombre de connexions per segon,
    \item nombre de bugs detectats en una versió,
    \item nombre de col·lisions en una taula hash.
  \end{itemize}
  \item Les distribucions discretes donen models simples per simular i raonar sobre aquests comptatges.
\end{itemize}
\column{0.39\textwidth}
\begin{block}{Idea clau}
Una variable aleatòria transforma un fenomen incert en un nombre que podem \important{calcular, simular i visualitzar}.
\end{block}
\vspace{0.2cm}
\begin{tikzpicture}
\node[draw=UABGreen,rounded corners,fill=UABLight,text width=4.5cm,align=center,inner sep=8pt] {
Sistema incert $\rightarrow$ variable aleatòria $X$ $\rightarrow$ distribució $\rightarrow$ decisió
};
\end{tikzpicture}
\end{columns}
\end{frame}

% ------------------------------------------------------------------
\begin{frame}{Objectius del tema}
Al final d'aquest tema hauríeu de poder:
\begin{enumerate}
  \item Definir una variable aleatòria discreta i la seva distribució.
  \item Calcular i interpretar esperança i variància.
  \item Reconèixer models de Bernoulli, binomial, geomètric i Poisson.
  \item Traduir problemes d'informàtica a variables aleatòries.
  \item Calcular probabilitats exactes amb fórmules i amb \R.
  \item Verificar resultats amb simulació Monte Carlo.
  \item Interpretar el paper de la variabilitat en sistemes computacionals.
\end{enumerate}
\end{frame}

% ------------------------------------------------------------------
\begin{frame}{Mapa del tema}
\begin{center}
\begin{tikzpicture}[node distance=1.15cm, every node/.style={font=\small}]
\node[draw,rounded corners,fill=UABLight,text width=3.0cm,align=center,minimum height=0.72cm] (va) {Variable aleatòria};
\node[draw,rounded corners,fill=SoftGray,below of=va,text width=3.0cm,align=center,minimum height=0.72cm] (dist) {Distribució discreta};
\node[draw,rounded corners,fill=SoftGray,below of=dist,text width=3.0cm,align=center,minimum height=0.72cm] (mom) {Esperança i variància};
\node[draw,rounded corners,fill=UABLight,right of=mom,xshift=3.2cm,text width=3.0cm,align=center,minimum height=0.72cm] (sim) {Simulació amb R};
\node[draw,rounded corners,fill=SoftGray,below of=mom,text width=2.5cm,align=center,minimum height=0.72cm] (bern) {Bernoulli};
\node[draw,rounded corners,fill=SoftGray,right of=bern,xshift=2.0cm,text width=2.5cm,align=center,minimum height=0.72cm] (bin) {Binomial};
\node[draw,rounded corners,fill=SoftGray,right of=bin,xshift=2.0cm,text width=2.5cm,align=center,minimum height=0.72cm] (geom) {Geomètrica};
\node[draw,rounded corners,fill=SoftGray,right of=geom,xshift=2.0cm,text width=2.5cm,align=center,minimum height=0.72cm] (pois) {Poisson};
\draw[->,thick,UABGreen] (va) -- (dist);
\draw[->,thick,UABGreen] (dist) -- (mom);
\draw[->,thick,AccentOrange] (mom) -- (sim);
\draw[->,thick,UABGreen] (mom) -- (bern);
\draw[->,thick,UABGreen] (mom) -- (bin);
\draw[->,thick,UABGreen] (mom) -- (geom);
\draw[->,thick,UABGreen] (mom) -- (pois);
\end{tikzpicture}
\end{center}
\begin{alertblock}{Pregunta transversal}
Quina distribució descriu millor el comptatge que observem?
\end{alertblock}
\end{frame}

% ------------------------------------------------------------------
\section{Variables aleatòries discretes}

\begin{frame}{Variable aleatòria: de l'espai mostral al nombre}
\begin{columns}[T]
\column{0.52\textwidth}
\begin{block}{Definició}
Una \important{variable aleatòria} és una funció
\[
X:\Omega\longrightarrow \mathbb{R}
\]
que assigna un valor numèric a cada resultat possible d'un experiment aleatori.
\end{block}
\begin{itemize}
  \item Els esdeveniments es poden escriure com $\{X=a\}$, $\{X\leq a\}$, etc.
  \item La distribució ens diu amb quina probabilitat pren cada valor.
\end{itemize}
\column{0.42\textwidth}
\begin{block}{Exemple clàssic}
Llancem dos daus i definim
\[
X=\text{suma dels resultats}.
\]
Aleshores $X$ pot valer $2,3,\ldots,12$.
\end{block}
\begin{block}{Exemple informàtic}
Enviem 100 paquets i definim
\[
X=\text{nombre de paquets perduts}.
\]
\end{block}
\end{columns}
\end{frame}

% ------------------------------------------------------------------
\begin{frame}{Distribució d'una variable discreta}
\begin{columns}[T]
\column{0.48\textwidth}
\begin{block}{Distribució}
Per a una variable discreta $X$:
\[
0\leq \Prob(X=a)\leq 1,
\qquad
\sum_a \Prob(X=a)=1.
\]
\end{block}
\begin{itemize}
  \item La distribució és una taula de valors i probabilitats.
  \item També es pot representar amb un gràfic de barres.
\end{itemize}
\column{0.46\textwidth}
\begin{center}
\begin{tabular}{c|cccccc}
\toprule
$x$ & 2 & 3 & 4 & $\cdots$ & 11 & 12 \\
\midrule
$\Prob(X=x)$ & $1/36$ & $2/36$ & $3/36$ & $\cdots$ & $2/36$ & $1/36$ \\
\bottomrule
\end{tabular}
\end{center}
\vspace{0.2cm}
\begin{tikzpicture}
\begin{axis}[
  width=6.2cm,height=3.1cm,
  ybar, bar width=7pt,
  ymin=0, ymax=0.18,
  xlabel={$x$}, ylabel={$\Prob(X=x)$},
  xtick={2,3,4,5,6,7,8,9,10,11,12},
  ytick={0,0.05,0.10,0.15},
  grid=major]
\addplot[fill=UABGreen!50,draw=UABGreen] coordinates {
(2,0.0278) (3,0.0556) (4,0.0833) (5,0.1111) (6,0.1389) (7,0.1667) (8,0.1389) (9,0.1111) (10,0.0833) (11,0.0556) (12,0.0278)};
\end{axis}
\end{tikzpicture}
\end{columns}
\end{frame}

% ------------------------------------------------------------------
\begin{frame}[fragile]{Distribucions discretes amb R}
\begin{lstlisting}
# Distribució exacta de la suma de dos daus
resultats <- expand.grid(dau1 = 1:6, dau2 = 1:6)
resultats$suma <- resultats$dau1 + resultats$dau2

table(resultats$suma)
prop.table(table(resultats$suma))
\end{lstlisting}

\vspace{0.2cm}
\begin{block}{Lectura computacional}
\begin{itemize}
  \item Generem tot l'espai mostral amb \texttt{expand.grid()}.
  \item Convertim resultats en valors de $X$.
  \item Comptem freqüències i les transformem en probabilitats.
\end{itemize}
\end{block}
\end{frame}

% ------------------------------------------------------------------
\begin{frame}{Esperança: valor mitjà esperat}
\begin{columns}[T]
\column{0.48\textwidth}
\begin{block}{Definició}
L'esperança d'una variable aleatòria discreta és
\[
\E(X)=\mu_X=\sum_a a\,\Prob(X=a).
\]
\end{block}
\begin{itemize}
  \item És el valor mitjà que esperaríem si repetíssim l'experiment moltes vegades.
  \item No ha de coincidir necessàriament amb un valor possible de $X$.
\end{itemize}
\column{0.46\textwidth}
\begin{block}{Interpretació informàtica}
Si $X$ és el nombre de tests que fallen en cada execució automàtica, $\E(X)$ és el nombre mitjà de fallades esperades per execució.
\end{block}
\begin{alertblock}{No és una predicció exacta}
És un resum del comportament mitjà a llarg termini.
\end{alertblock}
\end{columns}
\end{frame}

% ------------------------------------------------------------------
\begin{frame}{Variància: quant fluctua el resultat?}
\begin{columns}[T]
\column{0.48\textwidth}
\begin{block}{Definició}
La variància és
\[
\Var(X)=\sigma_X^2=\sum_a (a-\mu_X)^2\,\Prob(X=a).
\]
\end{block}
\begin{itemize}
  \item Mesura la dispersió respecte l'esperança.
  \item Una variància alta indica resultats més inestables.
\end{itemize}
\column{0.46\textwidth}
\begin{block}{Exemple computacional}
Dos algoritmes poden tenir el mateix temps mitjà però variàncies diferents:
\begin{itemize}
  \item un és estable,
  \item l'altre és ràpid sovint però molt lent en alguns casos.
\end{itemize}
\end{block}
\begin{alertblock}{Missatge}
La mitjana sola pot ser insuficient per a decidir.
\end{alertblock}
\end{columns}
\end{frame}

% ------------------------------------------------------------------
\begin{frame}{Mateixa esperança, risc diferent}
\begin{columns}[T]
\column{0.50\textwidth}
\begin{block}{Aposta A}
Guanyem $+1$ amb probabilitat $18/37$ i perdem $-1$ amb probabilitat $19/37$.
\[
\E(X)=\frac{18}{37}-\frac{19}{37}=-\frac{1}{37}.
\]
\end{block}
\column{0.44\textwidth}
\begin{block}{Aposta B}
Guanyem $+35$ amb probabilitat $1/37$ i perdem $-1$ amb probabilitat $36/37$.
\[
\E(Y)=35\frac{1}{37}-\frac{36}{37}=-\frac{1}{37}.
\]
\end{block}
\end{columns}
\vspace{0.2cm}
\begin{alertblock}{Conclusió}
Tenen la mateixa esperança, però no el mateix risc: $\Var(X)\approx 1.002$ i $\Var(Y)\approx 36.001$.
\end{alertblock}
\end{frame}

% ------------------------------------------------------------------
\begin{frame}[fragile]{Esperança i variància amb R}
\begin{lstlisting}
# Aposta A
x <- c(1, -1)
px <- c(18/37, 19/37)
EX <- sum(x * px)
VX <- sum((x - EX)^2 * px)

# Aposta B
y <- c(35, -1)
py <- c(1/37, 36/37)
EY <- sum(y * py)
VY <- sum((y - EY)^2 * py)

c(EX = EX, VarX = VX, EY = EY, VarY = VY)
\end{lstlisting}

\begin{block}{Missatge computacional}
Les fórmules són importants, però també ho és saber transformar-les en codi verificable.
\end{block}
\end{frame}

% ------------------------------------------------------------------
\section{Bernoulli i binomial}

\begin{frame}{Distribució de Bernoulli}
\begin{columns}[T]
\column{0.52\textwidth}
\begin{block}{Model}
Una variable Bernoulli només pot prendre dos valors:
\[
X=1 \quad \text{èxit}, \qquad X=0 \quad \text{fracàs}.
\]
\[
\Prob(X=1)=p, \qquad \Prob(X=0)=1-p=q.
\]
\[
\E(X)=p, \qquad \Var(X)=pq.
\]
\end{block}
\column{0.42\textwidth}
\begin{block}{Exemples informàtics}
\begin{itemize}
  \item Un test passa o falla.
  \item Un paquet arriba o es perd.
  \item Un usuari clica o no clica.
  \item Una petició HTTP retorna error o no.
\end{itemize}
\end{block}
\end{columns}
\end{frame}

% ------------------------------------------------------------------
\begin{frame}{Distribució binomial}
\begin{block}{Model}
Fem $n$ experiments Bernoulli independents, tots amb probabilitat d'èxit $p$.
Si $X$ és el nombre d'èxits, escrivim
\[
X\sim B(n,p).
\]
\[
\Prob(X=k)=\binom{n}{k}p^k(1-p)^{n-k},\qquad k=0,1,\ldots,n.
\]
\[
\E(X)=np,\qquad \Var(X)=np(1-p).
\]
\end{block}
\begin{alertblock}{Quan l'usem?}
Quan comptem quants èxits hi ha en un nombre fix d'intents independents.
\end{alertblock}
\end{frame}

% ------------------------------------------------------------------
\begin{frame}{Binomial en informàtica: paquets perduts}
\begin{columns}[T]
\column{0.52\textwidth}
\begin{block}{Plantejament}
Enviem $n=100$ paquets. Cada paquet es perd amb probabilitat $p=0.02$, independentment dels altres.
\[
X=\text{nombre de paquets perduts}\sim B(100,0.02).
\]
\end{block}
\begin{itemize}
  \item Nombre esperat de paquets perduts: $\E(X)=100\cdot 0.02=2$.
  \item Variància: $\Var(X)=100\cdot 0.02\cdot 0.98=1.96$.
\end{itemize}
\column{0.42\textwidth}
\begin{block}{Pregunta}
Quina és la probabilitat de perdre més de 5 paquets?
\[
\Prob(X>5)=1-\Prob(X\leq 5).
\]
\end{block}
\begin{alertblock}{Interpretació}
Això pot ajudar a fixar llindars d'alerta en una xarxa.
\end{alertblock}
\end{columns}
\end{frame}

% ------------------------------------------------------------------
\begin{frame}[fragile]{Binomial amb R: càlcul exacte i simulació}
\begin{lstlisting}
n <- 100
p <- 0.02

# Càlcul exacte
1 - pbinom(5, size = n, prob = p)

# Simulació Monte Carlo
set.seed(123)
B <- 100000
x <- rbinom(B, size = n, prob = p)
mean(x > 5)
mean(x)
var(x)
\end{lstlisting}

\begin{block}{Dues maneres complementàries}
\begin{itemize}
  \item \texttt{pbinom()} dona el càlcul exacte del model.
  \item \texttt{rbinom()} simula molts sistemes possibles.
\end{itemize}
\end{block}
\end{frame}

% ------------------------------------------------------------------
\begin{frame}{Distribució binomial: forma de la distribució}
\begin{center}
\begin{tikzpicture}
\begin{axis}[
  width=11.5cm,height=5.0cm,
  ybar, bar width=4pt,
  ymin=0, ymax=0.28,
  xlabel={$k$ paquets perduts}, ylabel={$\Prob(X=k)$},
  xtick={0,2,4,6,8,10},
  ytick={0,0.1,0.2},
  grid=major]
\addplot[fill=UABGreen!55,draw=UABGreen] coordinates {
(0,0.1326) (1,0.2707) (2,0.2734) (3,0.1823) (4,0.0902) (5,0.0353) (6,0.0114) (7,0.0031) (8,0.0007) (9,0.0001) (10,0.00002)};
\end{axis}
\end{tikzpicture}
\end{center}
\vspace{-0.1cm}
\begin{alertblock}{Lectura}
Encara que la mitjana sigui 2, poden aparèixer valors més grans. La distribució quantifica com de sorprenents són.
\end{alertblock}
\end{frame}

% ------------------------------------------------------------------
\section{Variable geomètrica}

\begin{frame}{Variable geomètrica}
\begin{columns}[T]
\column{0.52\textwidth}
\begin{block}{Model}
Fem experiments Bernoulli independents fins obtenir el primer èxit.
Si $X$ és el nombre de fracassos abans del primer èxit:
\[
\Prob(X=k)=q^k p,\qquad k=0,1,2,\ldots
\]
\[
\E(X)=\frac{q}{p},\qquad \Var(X)=\frac{q}{p^2}.
\]
\end{block}
\column{0.42\textwidth}
\begin{block}{Exemples informàtics}
\begin{itemize}
  \item Nombre de reintents abans que una connexió funcioni.
  \item Nombre de commits abans de detectar el primer test fallit.
  \item Nombre de peticions fins que apareix el primer error 500.
\end{itemize}
\end{block}
\end{columns}
\end{frame}

% ------------------------------------------------------------------
\begin{frame}{Geomètrica: reintents de connexió}
\begin{block}{Plantejament}
Una petició a un servei té probabilitat d'èxit $p=0.8$. Definim
\[
X=\text{nombre de fracassos abans del primer èxit}.
\]
Aleshores
\[
X\sim \operatorname{Geom}(p=0.8),\qquad \E(X)=\frac{0.2}{0.8}=0.25.
\]
\end{block}
\begin{columns}[T]
\column{0.48\textwidth}
\begin{itemize}
  \item $\Prob(X=0)=0.8$: funciona a la primera.
  \item $\Prob(X=1)=0.2\cdot0.8=0.16$.
  \item $\Prob(X\geq 3)=0.2^3=0.008$.
\end{itemize}
\column{0.46\textwidth}
\begin{alertblock}{Ús pràctic}
Serveix per raonar sobre polítiques de reintent, timeouts i degradació del servei.
\end{alertblock}
\end{columns}
\end{frame}

% ------------------------------------------------------------------
\begin{frame}[fragile]{Geomètrica amb R}
\begin{lstlisting}
p <- 0.8

# P(X = k), on X = fracassos abans del primer èxit
dgeom(0:5, prob = p)

# P(X >= 3)
1 - pgeom(2, prob = p)

# Simulació
set.seed(123)
x <- rgeom(100000, prob = p)
mean(x)
var(x)
\end{lstlisting}

\begin{block}{Detall important}
En \R, \texttt{rgeom()} segueix aquesta convenció: compta el nombre de fracassos abans del primer èxit.
\end{block}
\end{frame}

% ------------------------------------------------------------------
\section{Variable de Poisson}

\begin{frame}{Variable de Poisson}
\begin{columns}[T]
\column{0.52\textwidth}
\begin{block}{Model}
La distribució de Poisson modelitza el nombre d'esdeveniments en un interval de temps o espai:
\[
X\sim \operatorname{Pois}(\lambda).
\]
\[
\Prob(X=k)=\frac{\lambda^k e^{-\lambda}}{k!},\qquad k=0,1,2,\ldots
\]
\[
\E(X)=\lambda,\qquad \Var(X)=\lambda.
\]
\end{block}
\column{0.42\textwidth}
\begin{block}{Exemples informàtics}
\begin{itemize}
  \item Connexions per segon a un servidor.
  \item Errors per minut en un log.
  \item Incidències obertes per dia.
  \item Missatges que arriben a una cua.
\end{itemize}
\end{block}
\end{columns}
\end{frame}

% ------------------------------------------------------------------
\begin{frame}{Poisson: connexions a un servidor}
\begin{block}{Plantejament}
Un servidor rep de mitjana $2.5$ connexions per segon.
\[
X=\text{nombre de connexions en un segon}\sim \operatorname{Pois}(2.5).
\]
\end{block}
\begin{columns}[T]
\column{0.48\textwidth}
\[
\Prob(X=0)=e^{-2.5}\approx 0.082
\]
\[
\Prob(X=1)=2.5e^{-2.5}\approx 0.205
\]
\[
\Prob(X=2)=\frac{2.5^2e^{-2.5}}{2!}\approx 0.257
\]
\column{0.46\textwidth}
\begin{alertblock}{Pregunta operativa}
Quina és la probabilitat de rebre 3 o més connexions en un segon?
\[
\Prob(X\geq 3)=1-\Prob(X\leq 2).
\]
\end{alertblock}
\end{columns}
\end{frame}

% ------------------------------------------------------------------
\begin{frame}[fragile]{Poisson amb R}
\begin{lstlisting}
lambda <- 2.5

# Probabilitats puntuals
pois_vals <- dpois(0:8, lambda = lambda)
round(pois_vals, 3)

# Probabilitat de 3 o més connexions
1 - ppois(2, lambda = lambda)

# Si l'interval és de 2 segons, lambda es multiplica per 2
1 - dpois(0, lambda = 2 * lambda)
\end{lstlisting}

\begin{block}{Escalat temporal}
Si el ritme mitjà és $\lambda$ per segon, en $t$ segons el paràmetre és $t\lambda$.
\end{block}
\end{frame}

% ------------------------------------------------------------------
\begin{frame}{Distribució de Poisson: forma}
\begin{center}
\begin{tikzpicture}
\begin{axis}[
  width=11.2cm,height=4.8cm,
  ybar, bar width=8pt,
  ymin=0, ymax=0.28,
  xlabel={$k$ connexions en un segon}, ylabel={$\Prob(X=k)$},
  xtick={0,1,2,3,4,5,6,7,8,9,10},
  ytick={0,0.1,0.2},
  grid=major]
\addplot[fill=UABGreen!55,draw=UABGreen] coordinates {
(0,0.0821) (1,0.2052) (2,0.2565) (3,0.2138) (4,0.1336) (5,0.0668) (6,0.0278) (7,0.0099) (8,0.0031) (9,0.0009) (10,0.0002)};
\end{axis}
\end{tikzpicture}
\end{center}
\begin{alertblock}{Lectura}
La massa es concentra al voltant de $\lambda=2.5$, però hi ha probabilitat no nul·la de pics de càrrega.
\end{alertblock}
\end{frame}

% ------------------------------------------------------------------
\begin{frame}{Poisson i monitoratge: un llindar d'alerta}
\begin{columns}[T]
\column{0.54\textwidth}
\begin{block}{Context}
Suposem que en condicions normals hi ha $\lambda=2.5$ errors per minut.
Volem activar una alerta si observem un nombre molt alt d'errors.
\end{block}
\begin{itemize}
  \item Si el llindar és $X\geq 8$, calculem $\Prob(X\geq 8)$ sota el funcionament normal.
  \item Si aquesta probabilitat és petita, observar $8$ o més errors pot indicar una anomalia.
\end{itemize}
\column{0.40\textwidth}
\begin{alertblock}{Idea estadística}
Un llindar d'alerta és una decisió basada en una probabilitat de cua.
\end{alertblock}
\begin{block}{Amb R}
\[
\Prob(X\geq 8)=1-\Prob(X\leq 7).
\]
\texttt{1 - ppois(7, 2.5)}
\end{block}
\end{columns}
\end{frame}

% ------------------------------------------------------------------
\section{Simulació i pràctica computacional}

\begin{frame}[fragile]{Simulació Monte Carlo: quan el càlcul exacte no és còmode}
\begin{lstlisting}
set.seed(123)
B <- 100000

# Simulem 100000 segons d'activitat d'un servidor
connexions <- rpois(B, lambda = 2.5)

# Estimem probabilitats a partir de freqüències relatives
mean(connexions >= 3)
mean(connexions >= 8)

# Resums descriptius
mean(connexions)
var(connexions)
table(connexions)
\end{lstlisting}

\begin{block}{Connexió amb el tema anterior}
Les freqüències relatives de la simulació s'aproximen a les probabilitats del model.
\end{block}
\end{frame}

% ------------------------------------------------------------------
\begin{frame}[fragile]{Paral·lelització bàsica de simulacions amb R}
\begin{lstlisting}
library(future)
library(furrr)

plan(multisession, workers = 4)

simula_prob_cua <- function(i) {
  x <- rpois(100000, lambda = 2.5)
  mean(x >= 8)
}

set.seed(123)
resultats <- future_map_dbl(1:100, simula_prob_cua)

mean(resultats)
sd(resultats)
\end{lstlisting}

\begin{alertblock}{Per què és útil?}
Quan repetim moltes simulacions, bootstrap o experiments de rendiment, la computació paral·lela redueix el temps d'espera.
\end{alertblock}
\end{frame}

% ------------------------------------------------------------------
\begin{frame}{Quina distribució triem?}
\begin{center}
\begin{tabular}{p{3.2cm}p{4.1cm}p{4.4cm}}
\toprule
\textbf{Distribució} & \textbf{Pregunta típica} & \textbf{Exemple informàtic} \\
\midrule
Bernoulli & Èxit o fracàs? & Un test passa o falla \\
Binomial & Quants èxits en $n$ intents? & Paquets perduts en 100 enviaments \\
Geomètrica & Quants fracassos abans del primer èxit? & Reintents abans d'una connexió exitosa \\
Poisson & Quants esdeveniments en un interval? & Connexions per segon, errors per minut \\
\bottomrule
\end{tabular}
\end{center}
\vspace{0.2cm}
\begin{alertblock}{Estratègia}
Primer identifiquem què estem comptant; després escollim el model.
\end{alertblock}
\end{frame}

% ------------------------------------------------------------------
\begin{frame}{Mini-pràctica proposada}
\begin{block}{Objectiu}
Analitzar i simular un sistema amb comptatges aleatoris.
\end{block}
\begin{enumerate}
  \item Trieu un cas: paquets perduts, tests fallits, reintents de connexió o peticions a un servidor.
  \item Definiu la variable aleatòria i proposeu una distribució.
  \item Calculeu esperança, variància i una probabilitat rellevant.
  \item Simuleu el sistema amb \R.
  \item Compareu càlcul exacte i simulació.
  \item Genereu un gràfic i un informe breu amb Quarto.
  \item Pugeu el codi al repositori de \GitHub.
\end{enumerate}
\end{frame}

% ------------------------------------------------------------------
\begin{frame}{Tancament}
\begin{columns}[T]
\column{0.56\textwidth}
\begin{itemize}
  \item Les variables aleatòries discretes permeten modelitzar comptatges.
  \item L'esperança resumeix el comportament mitjà.
  \item La variància mesura la incertesa al voltant d'aquest comportament.
  \item Bernoulli, binomial, geomètrica i Poisson cobreixen molts casos bàsics.
  \item Amb \R podem calcular, simular i comunicar els resultats.
\end{itemize}
\column{0.38\textwidth}
\begin{block}{Missatge final}
La probabilitat no és només una fórmula: és una eina per dissenyar, monitorar i avaluar sistemes informàtics sota incertesa.
\end{block}
\vspace{0.2cm}
\begin{tikzpicture}
\node[draw=UABGreen,rounded corners,fill=UABLight,text width=4.4cm,align=center,inner sep=8pt] {
Model $+$ simulació $+$ visualització $=$ estadística computacional
};
\end{tikzpicture}
\end{columns}
\end{frame}

\end{document}
